\thispagestyle{empty}
\vspace*{2.5cm}
\begin{center}
  {\Large\bfseries Ringraziamenti}
\end{center}
\vspace{0.6cm}

\noindent Colgo questa pagina per fermare, brevemente, gli ultimi due anni --- e le persone
e i momenti che li hanno attraversati.

\medskip
\noindent Sono stati anni importanti, dentro e fuori l'università. Ne conservo immagini
spensierate: le giornate australiane passate a decidere cosa avremmo visitato, i pomeriggi
milanesi del dopo-esame a chiacchierare al bar. E frangenti in cui ho dovuto impegnarmi
davvero: le giornate lunghissime prima degli esami, passate a fare e rifare il giro del
programma, e i cinque esami in trenta giorni della prima sessione, che ricordo bene. Nel
complesso, di questo tempo porto con me un ricordo positivo, che mi ha arricchito.

\medskip
\noindent Tutto ciò che faccio, però, poggia sul sostegno della mia \textbf{famiglia}. Sono
sempre stato un ragazzo a cui la realtà di paese andava stretta, e ho sempre cercato di
andare oltre --- per curiosità, e per mettermi in gioco. La mia \textbf{famiglia} non mi ha
mai fatto mancare né i mezzi né il supporto per farlo. Non tutti possono dirlo: io sì, e
non lo do per scontato.

\medskip
\noindent Un grazie a tutte le persone che hanno fatto parte di questi due anni accademici:
i professori, i compagni, ma anche chi ha condiviso soltanto una conversazione nata per
caso, a Milano, ad Adelaide o ovunque sia capitato. Alcune di quelle conversazioni mi sono
rimaste impresse più di altre, ma tutte fanno parte del percorso, e lo hanno reso più ricco
e più interessante.

\medskip
\noindent Menzione d'onore, infine, per la mia cagnolina \textbf{Hope}, a cui voglio un
bene enorme: sono felice che sia stata accanto a me e alla mia \textbf{famiglia} in questi
anni. Ha reso tutto molto più bello.
